\documentclass[aspectratio=169]{beamer}
\usepackage[utf8]{inputenc}
\usepackage[T5]{fontenc}
\usepackage[vietnamese]{babel}
\usepackage{graphicx}
\usepackage{booktabs}
\usepackage{hyperref}
\usepackage{xcolor}

\usetheme{Madrid}
\usecolortheme{default}

\title[Tổng quan môn học]{ỨNG DỤNG TRÍ TUỆ NHÂN TẠO TRONG KẾ TOÁN \\ TỔNG QUAN MÔN HỌC}
\author{Giảng viên: \dots}
\institute{Đại học Đông Á}
\date{\today}

\begin{document}

\begin{frame}
    \titlepage
\end{frame}

\begin{frame}{Nội dung chính}
    \tableofcontents
\end{frame}

\section{Thông tin chung \& Đánh giá}
\begin{frame}{Mục tiêu môn học}
    \begin{itemize}
        \item Cung cấp kiến thức nền tảng về Trí tuệ nhân tạo (AI) và cách áp dụng trong Kế toán.
        \item Trang bị kỹ năng \textbf{Thực hành No-Code} với AI (ChatGPT, Copilot, Power Query, Advanced Data Analysis).
        \item Hiểu rõ và áp dụng quy trình phân tích dữ liệu MOSAIC.
        \item Nâng cao nhận thức về đạo đức, bảo mật và trách nhiệm pháp lý khi dùng AI.
    \end{itemize}
\end{frame}

\begin{frame}{Cấu trúc Điểm đánh giá}
    \begin{table}[]
        \centering
        \begin{tabular}{|l|c|p{7cm}|}
        \hline
        \textbf{Thành phần} & \textbf{Tỷ trọng} & \textbf{Hình thức} \\ \hline
        \textbf{Chuyên cần} & 10\% & Điểm danh, tham gia trên lớp \\ \hline
        \textbf{Thường xuyên} & 25\% & Bài trắc nghiệm (Quiz) và nộp Bài tập thực hành trên Canvas \\ \hline
        \textbf{Giữa kỳ} & 15\% & Kiểm tra thực hành / Tự luận \\ \hline
        \textbf{Cuối kỳ} & 50\% & Dự án nhóm (Project No-Code) - Thuyết trình và nộp báo cáo \\ \hline
        \end{tabular}
    \end{table}
\end{frame}

\section{Tài liệu học tập}
\begin{frame}{Tài liệu tham khảo chính}
    \textbf{1. Tài liệu Lý thuyết:}
    \begin{itemize}
        \item \textit{Artificial Intelligence in Accounting: Practical Applications} (Cory Ng) - Chương 1
        \item \textit{ChatGPT and AI for Accountants} (Scott Dell) - Chương 1, 2, 7, 10, 15
        \item \textit{Machine Learning in Accounting} (Ethan Blake) - Chương 4
        \item \textit{Accounting AI} (Charnelle Gibson) - Chương 5, 8, 9, 11
        \item \textit{Business Analytics: Solving Business Problems} (Arul Mishra) - Chương 11
        \item \textit{Artificial Intelligence for Audit, Forensic Accounting...} (Al Naqvi) - Chương 7, 8, 15
    \end{itemize}
\end{frame}

\begin{frame}{Tài liệu thực hành}
    \textbf{2. Tài liệu Thực hành \& Dữ liệu (No-Code):}
    \begin{itemize}
        \item \textit{Data and Analytics in Accounting} (Ann C. Dzuranin)
        \begin{itemize}
            \item Ch. 2: Kỹ năng phân tích dữ liệu nền tảng
            \item Ch. 7: Khám phá dữ liệu (Data Exploration)
            \item Ch. 9: Xây dựng Dashboard báo cáo tương tác
            \item Khung phân tích dữ liệu \textbf{MOSAIC} (Dự án cuối kỳ)
        \end{itemize}
        \item Bộ \textbf{Sổ tay Prompt Kế toán} (Tổng hợp từ giáo trình)
        \item Datasets mẫu: Sổ cái, Hóa đơn, Doanh thu (sử dụng trên Excel/Power Query).
    \end{itemize}
\end{frame}

\section{Lộ trình 15 buổi}

\begin{frame}{Lộ trình học tập - Phần 1: Nền tảng \& Nhập liệu}
    \begin{itemize}
        \item \textbf{Buổi 1: Tổng quan ứng dụng AI trong Kế toán}
        \begin{itemize}
            \item \textit{Thực hành:} Làm quen ChatGPT/Copilot. Viết prompt tóm tắt chuẩn mực.
        \end{itemize}
        \item \textbf{Buổi 2: Nhận diện \& Xử lý dữ liệu kế toán}
        \begin{itemize}
            \item \textit{Thực hành:} Làm sạch dữ liệu "bẩn" bằng Power Query (Excel) và ChatGPT.
        \end{itemize}
        \item \textbf{Buổi 3: AI hỗ trợ nhập liệu \& Xử lý chứng từ (OCR)}
        \begin{itemize}
            \item \textit{Thực hành:} Bóc tách dữ liệu hóa đơn (Hình ảnh/PDF) thành định dạng bảng.
        \end{itemize}
        \item \textbf{Buổi 4: AI xử lý công việc chi tiết (Định khoản)}
        \begin{itemize}
            \item \textit{Thực hành:} Đóng vai (Persona), định khoản nghiệp vụ tự động từ văn bản.
        \end{itemize}
        \vspace{0.3cm}
        \item \textbf{Đọc thêm (Sách Data \& Analytics):} Chương 1 (Tổng quan), Chương 3 \& 4 (Thu thập, chuẩn bị và làm sạch dữ liệu).
    \end{itemize}
\end{frame}

\begin{frame}{Lộ trình học tập - Phần 2: Phân tích \& Báo cáo}
    \begin{itemize}
        \item \textbf{Buổi 5: AI phát hiện sai sót và chẩn đoán bất thường}
        \begin{itemize}
            \item \textit{Thực hành:} Tìm điểm bất thường (ngày nghỉ, số chẵn) trong Sổ Nhật ký chung.
        \end{itemize}
        \item \textbf{Buổi 6: Ứng dụng AI phân tích Báo cáo Tài chính}
        \begin{itemize}
            \item \textit{Thực hành:} Tính tỷ số tài chính tự động và viết nhận xét "Sức khỏe tài chính".
        \end{itemize}
        \item \textbf{Buổi 7: AI hỗ trợ lập báo cáo theo yêu cầu (Dashboards)}
        \begin{itemize}
            \item \textit{Thực hành:} Lên outline Thuyết minh BCTC, vẽ Dashboard trực quan.
        \end{itemize}
        \item \textbf{Buổi 8: Ứng dụng AI trong Kế toán Quản trị}
        \begin{itemize}
            \item \textit{Thực hành:} Phân tích xu hướng \& dự báo dòng tiền/doanh thu các kịch bản.
        \end{itemize}
        \vspace{0.3cm}
        \item \textbf{Đọc thêm (Sách Data \& Analytics):} Chương 5 (Phân tích mô tả), Chương 6 (Phân tích chẩn đoán), Chương 8 (Phân tích dự báo).
    \end{itemize}
\end{frame}

\begin{frame}{Lộ trình học tập - Phần 3: Kiểm soát, Thuế \& Đạo đức}
    \begin{itemize}
        \item \textbf{Buổi 9: Ứng dụng AI trong Kiểm soát nội bộ}
        \begin{itemize}
            \item \textit{Thực hành:} Dò tìm "lỗ hổng" từ quy trình văn bản, đề xuất lưu đồ kiểm soát mới.
        \end{itemize}
        \item \textbf{Buổi 10: Ứng dụng AI trong Kiểm toán}
        \begin{itemize}
            \item \textit{Thực hành:} Lập "Bảng đánh giá rủi ro kiểm toán" bằng AI.
        \end{itemize}
        \item \textbf{Buổi 11: Ứng dụng AI trong Kế toán Thuế \& Tuân thủ}
        \begin{itemize}
            \item \textit{Thực hành:} Tóm tắt thông tư thuế dài, liệt kê điều kiện khấu trừ.
        \end{itemize}
        \item \textbf{Buổi 12: Đạo đức – Pháp lý – Trách nhiệm nghề nghiệp}
        \begin{itemize}
            \item \textit{Thực hành:} Phân tích tình huống đạo đức, kiểm tra tính "ảo giác" của AI.
        \end{itemize}
        \vspace{0.3cm}
        \item \textbf{Đọc thêm (Sách Data \& Analytics):} Các chương còn lại (Ứng dụng Analytics trong Kiểm toán, Thuế \& Đạo đức nghề nghiệp).
    \end{itemize}
\end{frame}

\section{Dự án cuối kỳ}
\begin{frame}{Dự án cuối kỳ (Buổi 13 - 15)}
    \textbf{Project No-Code: Vận dụng toàn diện AI vào Dữ liệu Kế toán}
    \begin{itemize}
        \item \textbf{Dữ liệu thực tế:} Mỗi nhóm nhận một Raw Dataset (Sổ cái, Hóa đơn).
        \item \textbf{Quy trình MOSAIC:}
        \begin{enumerate}
            \item \textbf{Motivation \& Objective:} Xác định mục tiêu phân tích.
            \item \textbf{Strategy:} Chọn chiến lược xử lý dữ liệu với AI.
            \item \textbf{Analyze:} Dùng AI làm sạch, định khoản và tính toán.
            \item \textbf{Interpret:} Phân tích kết quả, lập báo cáo.
            \item \textbf{Communicate:} Xây dựng Dashboard và Thuyết trình (Buổi 14-15).
        \end{enumerate}
        \item \textbf{Yêu cầu:} Vận dụng kỹ thuật Prompt Engineering \& Advanced Data Analysis đã học.
    \end{itemize}
\end{frame}

\begin{frame}{Tổng kết}
    \begin{center}
        \Large \textbf{Chào mừng các bạn đến với môn học!} \\
        \vspace{0.5cm}
        \normalsize Cùng sẵn sàng bứt phá với Kế toán trong kỷ nguyên AI.
    \end{center}
\end{frame}

\end{document}
